\uline{flat morphism} @@M0@@ \uline{such that every irreducible component} @@M1@@ \uline{of} @@M2@@\uline{, endowed with the induced reduced structure, is integral over} @@M3@@.

Indeed, there is a finite extension @@M4@@ of the function field @@M5@@ such that every irreducible component of @@M6@@ is geometrically integral (by EGA IV 4.6.8). Let @@M7@@ be an integral @@M8@@-scheme such that @@M9@@ and such that there are a nonempty open subset @@M10@@ of @@M11@@ and a finite flat morphism @@M12@@. There is a nonempty open subset @@M13@@ of @@M14@@ such that every irreducible component of @@M15@@, endowed with the induced reduced structure, is integral over @@M16@@ (EGA IV 9.7.7). Take @@M17@@ and @@M18@@.

Let us now prove 1.2. Let @@M19@@ and @@M20@@ be as in 1.3. For every @@M21@@, let @@M22@@ be a Chow presentation with @@M23@@ integral and @@M24@@ birational. After shrinking @@M25@@ and @@M26@@, we may suppose that @@M27@@ is integral over @@M28@@ (EGA IV 9.7.7). Let @@M29@@ be the canonical morphism. Then @@M30@@ is representable by a group scheme that is a sum of quasi-projective open subschemes of finite type.

Finally, 1.1 says that, after shrinking @@M31@@ and @@M32@@, the morphism @@M33@@ is representable by a quasi-affine morphism. Consequently, @@M34@@ is representable by a group scheme that is a sum of quasi-projective open subschemes of finite type; by finite flat descent (SGA 1 VIII 7.6 and EGA II 6.6.1), the same holds for @@M35@@.

\sgaStableTarget{sga6.e12.lem.1-4}\textbf{Lemma 1.4.} \uline{Let} @@M36@@ \uline{be a field,} @@M37@@ \uline{and} @@M38@@ \uline{two} @@M39@@\uline{-group schemes locally of finite type, and let} @@M40@@ \uline{be a quasi-affine morphism. Then} @@M41@@ \uline{is affine.}

Indeed, let @@M42@@; @@M43@@ is a @@M44@@-group scheme of finite type that is quasi-affine; by (SGA 3 @@M45@@ 11.11), @@M46@@ is therefore affine.

